%!TEX root = bertini_real_manual.tex

\subsection{Instructions for OSX}

These instructions were last updated in July 2025.


\subsubsection{Preliminary Work}
\begin{itemize}
  \item Compilers \&c: In terminal, type {\tt xcode-select --install}
   
  \item Homebrew: It is highly recommended that Mac users install the program \href{http://brew.sh}{Homebrew} to use to install these packages. Once that has been done, installing the previously listed dependencies becomes simple. In terminal,  type \texttt{brew search \_\_\_\_} to list packages related to \_\_\_\_, where \_\_\_\_ is your search (for example, GMP, Boost, or MPI). To download new software via Homebrew, type in terminal \texttt{brew install \_\_\_\_} 

  Suggested call:

  \texttt{brew install cmake autoconf automake libtool mpfr gmp mpich boost}
   \item Matlab: make sure to have it installed, and on the path for the command line.
   \item Python: use {\tt pip} to install dependencies
\end{itemize}


\subsubsection{Build and install Bertini 1}

\begin{enumerate}
  \item Download the source tarball ({\tt .tar.gz}) from \href{http://bertini.nd.edu}{bertini.nd.edu}
  \item Unpack the tarball (just double click it, if you're into mousing)
  \item If you didn't already, now use Homebrew to install the dependencies for Bertini1 \\{\tt brew install cmake autoconf automake libtool mpfr gmp mpich boost}\\ (Bertini1 doesn't depend on Boost, but Bertini\_real does)
  \item In the terminal, move to the unpacked folder for Bertini 1

  \item Compile.
  \begin{enumerate}

    \item For Bertini versions 1.6.x and before, use the make build system, using this classic pattern:
    \begin{enumerate}
      \item {\tt ./configure}.  \\ If you're using homebrew, this may require \\ \verb|./configure LDFLAGS=-L/opt/homebrew/lib CPPFLAGS=-I/opt/homebrew/include| 
      \item {\tt make -j 4 \&\& make install}\\ installing may require \verb|sudo|
    \end{enumerate}

    \item For Bertini versions 1.7 and later, use CMake.
      \begin{enumerate}
        \item \verb|mkdir build && cd build|
        \item \verb|cmake ../|
        \item \verb|make -j 4|
        \item \verb|make install| \\ Installing may require \verb|sudo|
      \end{enumerate}
  \end{enumerate}
\end{enumerate}

\subsubsection{Build and install Bertini\_real}

\begin{enumerate}
  \item Clone the repo to get the code.
  \begin{enumerate}
    \item {\tt which git}  if it's emtpy, then {\tt brew install git}
    \item {\tt mkdir code \&\& cd code}
    \item {\tt git clone https://github.com/ofloveandhate/bertini\_real}
  \end{enumerate}

  \item Move into the directory.  {\tt cd bertini\_real}

  \item Compile using one of two options
    \begin{enumerate}
      \item For Bertini\_real 1.8 and before, using the classic make system:
      \begin{enumerate}
        \item {\tt autoreconf -i}
        \item {\tt ./configure}
        \item {\tt make -j 4 \&\& make install} Installing may require \verb|sudo|
      \end{enumerate} 

      \item For Bertini\_real 1.9 and later, using cmake:
      \begin{enumerate}
        \item \verb|mkdir build && cd build|
        \item \verb|cmake ../|
        \item \verb|make -j 4|
        \item \verb|make install| \\ Installing may require \verb|sudo|
      \end{enumerate}


    \end{enumerate}

\end{enumerate}


Problems?  Raise an issue \href{https://github.com/ofloveandhate/bertini_real/issues}{on Github}, please. 
